\begin{abstract}
Neural predictive process monitoring (PPM) models are compared almost only
through task metrics such as next-activity accuracy or suffix similarity. Yet
every such model maps each running prefix to a latent vector, so that a case
traces a trajectory through representation space. What these representations
contain has not been examined. We ask three questions. What do PPM models
learn in their prefix representations? Does the geometry of these
representations carry information about prediction errors beyond confidence?
Do models place prefixes with similar futures close together even when their
histories differ? To answer them we introduce a synthetically validated
measurement framework. It combines collapse diagnostics, trajectory metrics,
nested reliability models and history-controlled tests of \emph{future
equivalence}. We apply it to nine configurations on five public event logs:
seven published next-event and suffix models reproduced under one
chronological split, and two controlled Transformers that differ only in
their objective. The study connects machine learning for process mining with
representation analysis, and it targets two practical needs: knowing when a
prediction can be trusted, and knowing what a model's state says about the
future of a case.
\TODO{Add the main results (one to two sentences with the key numbers) once
the analyses are complete, as the CFP requires the abstract to state the
results achieved.}
\keywords{Predictive process monitoring \and Representation geometry \and
Prediction reliability \and Suffix prediction \and Future equivalence}
\end{abstract}
